\chapter{Evaluation}
\label{cha:evaluation}

This chapter evaluates the prototype described in Chapter~\ref{cha:impl}. The evaluation targets
the third objective of this dissertation (Section~\ref{sec:int_objectives}): to measure the
overhead introduced by the integration and to assess the correctness of the resolution mechanism.
Because the contribution of this work is the \emph{unification} layer (the function registry and
the deployment-time resolution cascade) and the AWS provider added to QuickFaaS, the evaluation
concentrates on the cost of resolving internal-function endpoints, on how that cost scales, and on
the size overhead of the AWS agnostic packaging layer.

\section{Scope and Goals}
\label{sec:eval-scope}

The measurements exercise only \emph{local} operations: constructing a workflow model, traversing
it, resolving internal calls against the registry, reading and writing the registry file, and
packaging a Lambda bundle. Operations that call the cloud provider (Lambda/IAM/S3 requests, Cloud
Run/Cloud Functions APIs, the QuickFaaS subprocess) are deliberately excluded, for two reasons:
they are dominated by network and provider-side latency measured in seconds, and they are
non-deterministic and out of the unit-test scope by design. The evaluation therefore answers a
focused question: \emph{how much does the unification layer add to a deployment, and how does that
overhead scale with the number of calls and the size of the registry?}

Three families of experiments are reported:
\begin{itemize}
  \item \textbf{Resolution overhead and scalability} --- the cost of resolving internal calls as a
  function of the number of calls ($N$) and the size of the registry ($R$), and the effect of the
  single-read optimization.
  \item \textbf{Real auto-deploy resolvers} --- the cost of the production resolvers
  (\texttt{Aws\allowbreak Internal\allowbreak Function\allowbreak Resolver}, \texttt{Workflow\allowbreak Internal\allowbreak Function\allowbreak Resolver}) and of the
  registry write path.
  \item \textbf{Packaging overhead} --- the size the AWS provider-agnostic layer adds to a Lambda
  deployment bundle.
\end{itemize}

\section{Methodology}
\label{sec:eval-methodology}

Benchmarks were implemented with the Java Microbenchmark Harness (JMH) 1.37, using
\texttt{AverageTime} mode with results reported in microseconds per operation ($\mu$s/op). Each
measured value is consumed through a \texttt{Blackhole} to prevent dead-code elimination by the
JIT compiler. Unless otherwise stated, runs used one JVM fork with three warm-up iterations and
five measurement iterations of one second each (\texttt{-f 1 -wi 3 -i 5 -w 1 -r 1}), so that the
already-JIT-compiled code is measured. Selected write-path experiments (P13, P17) used three forks
(\texttt{-f 3}) to reduce noise.

\paragraph{Caveat on absolute values.}
Measurements were taken on a shared, non-dedicated machine with a deliberately short
configuration. The absolute numbers are therefore meaningful for comparing \emph{trends and
relationships} (linear vs.\ quadratic growth, before vs.\ after an optimization), not as
hardware-definitive figures. Thesis-grade absolute numbers would require repeating the suite with
\texttt{-f 3} on a dedicated machine.

\paragraph{Notation.}
The experiments share five symbols; what changes between them is the \emph{relationship} among
these quantities, summarised in Table~\ref{tab:eval-notation}.

\begin{table}[htbp]
  \centering
  \caption{Notation used throughout the evaluation.}
  \label{tab:eval-notation}
  \begin{tabular}{cl}
    \hline
    \textbf{Symbol} & \textbf{Meaning} \\
    \hline
    $N$ & total number of \texttt{call} steps in the workflow (internal or external), $N = I + E$ \\
    $R$ & number of entries in the registry file when it is read/written \\
    $F$ & number of \emph{distinct} internal functions the workflow references \\
    $I$ & number of internal calls (repetitions counted) \\
    $E$ & number of external calls (repetitions counted) \\
    \hline
  \end{tabular}
\end{table}

$R$ is the physical size of the registry file and is not the same as $F$: a registry may hold
$R{=}1000$ entries while a workflow references only $F{=}10$ of them. Internal calls are
distributed round-robin over the $F$ distinct functions.

\section{Resolution Overhead and Scalability}
\label{sec:eval-resolution}

\paragraph{Cost of the unification (P3).}
The first experiment isolates the cost of resolving an internal call against the cost of an
external call (which bypasses the registry). Resolving an internal call costs about
$7.5\,\mu$s and scales linearly with the number of calls, because each call re-reads and re-parses
the entire registry file from disk; an external call costs about $0.01\,\mu$s
(Table~\ref{tab:eval-p3}). In the general case the cost grows in direct proportion to both the
number of calls and the number of registry entries.

\begin{table}[htbp]
  \centering
  \caption{P3 --- resolution cost of internal vs.\ external calls, with $R{=}1$.}
  \label{tab:eval-p3}
  \begin{tabular}{rrr}
    \hline
    \textbf{Calls ($N$)} & \textbf{Internal --- resolved ($\mu$s)} & \textbf{External --- not resolved ($\mu$s)} \\
    \hline
    1   & 7.5    & 0.02 \\
    10  & 74.6   & 0.14 \\
    50  & 368.0  & 0.65 \\
    100 & 747.2  & 1.10 \\
    200 & 1511.4 & 2.34 \\
    \hline
  \end{tabular}
\end{table}

\paragraph{Effect of registry size (P6).}
Fixing the workflow and varying the registry size shows that the per-call cost grows with $R$ and
is essentially independent of $N$: about $180\,\mu$s of fixed I/O and parsing per read plus about
$0.56\,\mu$s per registered entry (Table~\ref{tab:eval-p6}). The four $N$ curves collapse onto a
single line, confirming that the cost grows in direct proportion to both the number of calls and
the registry size, where the fixed cost of reopening the file dominates for realistic registry
sizes.

\begin{table}[htbp]
  \centering
  \caption{P6 --- per-call resolution cost as a function of registry size $R$ (independent of $N$).}
  \label{tab:eval-p6}
  \begin{tabular}{rr}
    \hline
    \textbf{Registry size ($R$)} & \textbf{Cost per call ($\mu$s)} \\
    \hline
    1    & 181.0 \\
    10   & 185.1 \\
    50   & 207.0 \\
    200  & 279.2 \\
    1000 & 741.7 \\
    \hline
  \end{tabular}
\end{table}

\paragraph{The single-read optimization (P7--P9).}
Reading the registry once per workflow into an in-memory snapshot and resolving against it
eliminates that call-times-registry-size product: the cost no longer multiplies the two factors
together, becoming instead one read proportional to the registry size plus a cheap constant-time
lookup per call. Table~\ref{tab:eval-p7} compares, in the same run, resolving
the registry per call (naive) against a single read per resolution (optimized). With the snapshot,
the cost becomes almost independent of $N$, yielding a speedup of up to $\sim$200$\times$ in the
worst corner ($N{=}200$, $R{=}1000$: $\sim$105\,ms down to $\sim$0.54\,ms). Experiments P8 and P9
confirm the same trend on real \texttt{Workflow} objects rather than the isolated store primitive.

\begin{table}[htbp]
  \centering
  \caption{P7 --- total resolution time, naive (per-call read) vs.\ optimized (single read).}
  \label{tab:eval-p7}
  \begin{tabular}{lrrrr}
    \hline
     & $N{=}1$ & $N{=}10$ & $N{=}50$ & $N{=}200$ \\
    \hline
    naive, $R{=}1$       & 4.6   & 47.0   & 233.6   & 920.2    \\
    optimized, $R{=}1$   & 4.7   & 5.0    & 5.0     & 5.7      \\
    naive, $R{=}1000$    & 543.7 & 5480.8 & 26698.5 & 104889.0 \\
    optimized, $R{=}1000$& 557.4 & 535.7  & 551.8   & 536.6    \\
    \hline
  \end{tabular}
\end{table}

\paragraph{Confirming the optimization on real workflows (P8).}
P8 exercises the resolution of a real \texttt{Workflow} object, not the isolated store primitive,
across a grid of $1$ to $50$ distinct functions and $1$ to $200$ calls, with the registry always
hitting ($R{=}F$). The naive path stays dominated by the number of calls and is largely
indifferent to how many distinct functions are involved, at roughly $200\,\mu$s per call at low
$N$ (Table~\ref{tab:eval-p8-naive}), while the optimized path stays close to flat
(Table~\ref{tab:eval-p8-optimized}); the measured speedup climbs from about $1\times$ at $N{=}1$
to about $110\times$ at $N{=}200$.

\begin{table}[htbp]
  \centering
  \caption{P8 --- naive resolution time ($\mu$s) of a real workflow, $F$ distinct functions vs.\
  $N$ calls, $R{=}F$.}
  \label{tab:eval-p8-naive}
  \begin{tabular}{rrrrrr}
    \hline
    $F\backslash N$ & 1 & 5 & 10 & 50 & 200 \\
    \hline
    1  & 216 & 1\,118 & 2\,103 & 9\,544  & 38\,394 \\
    2  & 192 & 956    & 1\,902 & 9\,555  & 37\,950 \\
    5  & 189 & 995    & 1\,911 & 9\,663  & 38\,468 \\
    10 & 195 & 1\,132 & 2\,390 & 10\,544 & 39\,340 \\
    20 & 199 & 997    & 2\,096 & 10\,017 & 41\,205 \\
    50 & 219 & 1\,133 & 2\,194 & 10\,857 & 43\,360 \\
    \hline
  \end{tabular}
\end{table}

\begin{table}[htbp]
  \centering
  \caption{P8 --- optimized (single-read) resolution time ($\mu$s) of a real workflow, same grid
  as Table~\ref{tab:eval-p8-naive}.}
  \label{tab:eval-p8-optimized}
  \begin{tabular}{rrrrrr}
    \hline
    $F\backslash N$ & 1 & 5 & 10 & 50 & 200 \\
    \hline
    1  & 191 & 195 & 197 & 228 & 337 \\
    2  & 191 & 193 & 200 & 230 & 341 \\
    5  & 191 & 194 & 209 & 228 & 349 \\
    10 & 228 & 197 & 202 & 241 & 362 \\
    20 & 199 & 202 & 215 & 237 & 349 \\
    50 & 224 & 221 & 223 & 251 & 367 \\
    \hline
  \end{tabular}
\end{table}

On a realistic workflow ($F{=}10$, $N{=}50$ fixed), P9 isolates the effect of the registry size:
the naive cost climbs with $R$ while the optimized cost stays almost flat, so the speedup grows
from $\sim$41$\times$ at $R{=}10$ to $\sim$58$\times$ at $R{=}1000$ (Table~\ref{tab:eval-p9}).

\begin{table}[htbp]
  \centering
  \caption{P9 --- resolution time ($\mu$s) vs.\ registry size $R$, $F{=}10$/$N{=}50$ fixed.}
  \label{tab:eval-p9}
  \begin{tabular}{rrrr}
    \hline
    \textbf{$R$} & \textbf{naive} & \textbf{optimized} & \textbf{speedup} \\
    \hline
    10   & 9\,829  & 237 & 41.4 \\
    50   & 10\,899 & 268 & 40.7 \\
    200  & 14\,964 & 324 & 46.2 \\
    1000 & 43\,245 & 749 & 57.8 \\
    \hline
  \end{tabular}
\end{table}

\section{Real Auto-Deploy Resolvers and Write Path}
\label{sec:eval-real-resolvers}

\paragraph{Production resolvers (P10, P11).}
Experiments P10 and P11 measure the resolvers that actually run during deployment, one per
provider: \texttt{Aws\allowbreak Internal\allowbreak Function\allowbreak Resolver} (Table~\ref{tab:eval-p10}) and
\texttt{Workflow\allowbreak Internal\allowbreak Function\allowbreak Resolver} (Table~\ref{tab:eval-p11}), with the
registry always hitting ($R{=}F$). Both now carry the single-read optimization demonstrated in
P7--P9 (\S\ref{sec:eval-resolution}): each resolver reads the registry once into a snapshot per
\texttt{resolve()} call and looks up every internal call against it
(\texttt{FunctionRegistryStore.tryResolveEntryIn}) instead of re-reading the file per call. The
result is the same shape as the already-optimized reference implementation: cost is dominated by
$N$ and roughly flat in $F$, at magnitudes comparable to P8's optimized path rather than the
naive, call-times-registry-size cost the unoptimized versions of these classes used to pay. The
two providers land in the same broad range (a few tens to a few thousand microseconds depending on
$N$), with no consistent directional gap between them once the registry-read cost is amortized
--- what asymmetry remains is within the run-to-run noise of this evaluation's shared,
non-dedicated machine (\S\ref{sec:eval-methodology}), most visible at $F{=}10,N{=}5$ in
Table~\ref{tab:eval-p10} and $F{=}50,N{=}50$, both flagged by their own unusually wide error
margins rather than reflecting a code effect.

\begin{table}[htbp]
  \centering
  \caption{P10 --- AWS provider: resolution time ($\mu$s) of the real resolver
  (\texttt{Aws\allowbreak Internal\allowbreak Function\allowbreak Resolver}), $R{=}F$, after the
  single-read optimization was propagated into it.}
  \label{tab:eval-p10}
  \begin{tabular}{lrrrrr}
    \hline
    $F\backslash N$ & 1 & 5 & 10 & 50 & 200 \\
    \hline
    1  & 17 & 44  & 124 & 1\,240 & 3\,486 \\
    10 & 19 & 171 & 85  & 424    & 1\,702 \\
    50 & 52 & 48  & 98  & 1\,669 & 1\,889 \\
    \hline
  \end{tabular}
\end{table}

\begin{table}[htbp]
  \centering
  \caption{P11 --- GCP provider: resolution time ($\mu$s) of the real resolver
  (\texttt{Workflow\allowbreak Internal\allowbreak Function\allowbreak Resolver}), $R{=}F$, after
  the single-read optimization was propagated into it.}
  \label{tab:eval-p11}
  \begin{tabular}{lrrrrr}
    \hline
    $F\backslash N$ & 1 & 5 & 10 & 50 & 200 \\
    \hline
    1  & 20 & 47  & 83  & 353 & 1\,417 \\
    10 & 47 & 105 & 213 & 963 & 2\,966 \\
    50 & 31 & 61  & 99  & 402 & 1\,573 \\
    \hline
  \end{tabular}
\end{table}

\paragraph{Registry write path (P13, P17).}
Writing is more expensive than resolving. Every \texttt{put()} re-reads and rewrites the entire
registry file, so writing $K$ functions into a registry that already holds $R_0$ entries takes
longer than linear time in $K$ (Table~\ref{tab:eval-p13}): each successive write pays for a
registry that has grown by one entry since the write before it, so the cost compounds as $K$
grows, on top of the floor set by $R_0$. The miss-then-write pattern the real resolvers perform
before any cloud call (\texttt{tryResolveEntry} followed by \texttt{put}) adds a further
24--54\% on top (Table~\ref{tab:eval-p17}), since \texttt{tryResolveEntry} is itself a
registry-size-proportional scan. In absolute terms these values remain
in the hundreds of microseconds to low hundreds of milliseconds even for $K{=}100$ successive
writes. The $R_0{=}1000$ column in Table~\ref{tab:eval-p13} does not follow the otherwise
monotonic pattern at $K{=}100$; this single data point is attributed to noise on the shared,
non-dedicated benchmarking machine (\S\ref{sec:eval-methodology}), not to a code effect.

\begin{table}[htbp]
  \centering
  \caption{P13 --- total write time ($\mu$s): $K$ successive \texttt{put} calls, initial registry
  size $R_0$.}
  \label{tab:eval-p13}
  \begin{tabular}{lrrrrr}
    \hline
    $K\backslash R_0$ & 0 & 10 & 50 & 200 & 1000 \\
    \hline
    1   & 2\,673   & 2\,860   & 3\,235   & 4\,139   & 2\,761   \\
    10  & 27\,691  & 26\,986  & 29\,912  & 39\,164  & 24\,029  \\
    100 & 329\,757 & 304\,187 & 330\,100 & 419\,531 & 237\,910 \\
    \hline
  \end{tabular}
\end{table}

\begin{table}[htbp]
  \centering
  \caption{P17 --- total write time ($\mu$s): $K$ successive miss+\texttt{put} pairs, initial
  registry size $R_0$.}
  \label{tab:eval-p17}
  \begin{tabular}{lrrrrr}
    \hline
    $K\backslash R_0$ & 0 & 10 & 50 & 200 & 1000 \\
    \hline
    1   & 3\,309   & 3\,545   & 3\,959   & 5\,122   & 4\,243   \\
    10  & 35\,107  & 34\,358  & 37\,190  & 47\,631  & 47\,817  \\
    100 & 410\,389 & 382\,140 & 403\,488 & 511\,878 & 346\,068 \\
    \hline
  \end{tabular}
\end{table}

\section{Match Strategy, Mixed Workflows, and Structure}
\label{sec:eval-secondary}

\paragraph{Exact vs.\ suffix match (P14).}
The suffix-match path (regional disambiguation) fails the exact lookup and scans the loaded map,
reintroducing a cost proportional to the registry size, but over an already-loaded map: the ratio
to exact match grows to only $\sim$1.9$\times$ at $R{=}1000$ and remains cheap in absolute terms
(Table~\ref{tab:eval-p14}).

\begin{table}[htbp]
  \centering
  \caption{P14 --- exact vs.\ suffix match ($\mu$s), $F{=}10$/$N{=}50$ fixed.}
  \label{tab:eval-p14}
  \begin{tabular}{rrrr}
    \hline
    \textbf{$R$} & \textbf{exact} & \textbf{suffix} & \textbf{ratio} \\
    \hline
    10   & 219 & 228   & 1.04 \\
    100  & 259 & 333   & 1.29 \\
    200  & 309 & 442   & 1.43 \\
    1000 & 768 & 1\,456 & 1.90 \\
    \hline
  \end{tabular}
\end{table}

\paragraph{Mixed internal/external workflows (P15, P16).}
The first workflows to mix internal and external calls confirm that resolution cost scales with
$I$ (internal calls), not with $N{=}I{+}E$: the $I{=}0$ line stays flat as $E$ grows to 200,
because external calls incur no registry lookup (Table~\ref{tab:eval-p15}). The registry-size
effect proved in P9 is unchanged, within the same order of magnitude, when a fraction of the
calls are external (Table~\ref{tab:eval-p16}).

\begin{table}[htbp]
  \centering
  \caption{P15 --- resolution time ($\mu$s) of a mixed workflow, $R{=}F{=}10$ fixed;
  $I$ internal vs.\ $E$ external calls.}
  \label{tab:eval-p15}
  \begin{tabular}{rrrrrr}
    \hline
    $I\backslash E$ & 0 & 1 & 10 & 50 & 200 \\
    \hline
    0   & 181 & 190 & 187 & 183 & 183 \\
    10  & 189 & 205 & 203 & 189 & 255 \\
    50  & 217 & 217 & 230 & 219 & 322 \\
    200 & 326 & 321 & 323 & 315 & 393 \\
    \hline
  \end{tabular}
\end{table}

\begin{table}[htbp]
  \centering
  \caption{P16 --- resolution time ($\mu$s) vs.\ registry size $R$, mixed workflow with
  $I{=}40$/$E{=}10$/$F{=}10$ fixed.}
  \label{tab:eval-p16}
  \begin{tabular}{rr}
    \hline
    \textbf{$R$} & \textbf{Time ($\mu$s)} \\
    \hline
    10   & 346 \\
    50   & 247 \\
    200  & 306 \\
    1000 & 776 \\
    \hline
  \end{tabular}
\end{table}

\paragraph{Workflow shape (P18, P19).}
Nesting depth (P18) leaves the time essentially constant ($\sim$200\,$\mu$s across depths 0--5),
and parallel width (P19) grows only with the total number of internal leaf calls
($N = 5\times\text{width}$), in the same order of magnitude as the flat workflows of P8/P15
(Tables~\ref{tab:eval-p18} and~\ref{tab:eval-p19}). Although the resolver recurses explicitly into
\texttt{Iteration} and \texttt{Parallel} contexts, the shape of the tree is irrelevant: only the
count of internal calls matters.

\begin{table}[htbp]
  \centering
  \caption{P18 --- resolution time ($\mu$s) vs.\ nesting depth, 20 internal calls fixed,
  $R{=}F{=}10$.}
  \label{tab:eval-p18}
  \begin{tabular}{rr}
    \hline
    \textbf{Depth} & \textbf{Time ($\mu$s)} \\
    \hline
    0 & 207.3 \\
    2 & 200.7 \\
    3 & 200.4 \\
    5 & 199.6 \\
    \hline
  \end{tabular}
\end{table}

\begin{table}[htbp]
  \centering
  \caption{P19 --- resolution time ($\mu$s) vs.\ \texttt{Parallel} branch width,
  $N = 5\times$ width, $R{=}F{=}10$.}
  \label{tab:eval-p19}
  \begin{tabular}{rrr}
    \hline
    \textbf{Branches} & \textbf{$N$} & \textbf{Time ($\mu$s)} \\
    \hline
    1   & 5   & 188.3 \\
    10  & 50  & 220.8 \\
    20  & 100 & 276.2 \\
    100 & 500 & 583.7 \\
    \hline
  \end{tabular}
\end{table}

\section{Packaging Overhead (S1)}
\label{sec:eval-bundle}

The AWS provider added in this work packages a Lambda through the QuickFaaS agnostic layer (the
developer hook plus the \texttt{AwsHttpTemplate} adapter and \texttt{aws-lambda-java-core}). Table
\ref{tab:eval-s1} compares the resulting bundle against an equivalent native \texttt{RequestHandler}
implementation of the same logic. The agnostic layer adds about $0.9$\,KB --- 6.6\% of a trivial
bundle, and well under 0.1\% of a real function --- consistent with the original QuickFaaS
evaluation for GCP and Azure.

\begin{table}[htbp]
  \centering
  \caption{S1 --- Lambda bundle size, AWS agnostic layer vs.\ native implementation.}
  \label{tab:eval-s1}
  \begin{tabular}{lr}
    \hline
    \textbf{Bundle} & \textbf{Size} \\
    \hline
    AWS agnostic (QuickFaaS hook + \texttt{AwsHttpTemplate} + \texttt{aws-lambda-java-core}) & 14\,552 bytes \\
    AWS native (single \texttt{RequestHandler}, same logic) & 13\,647 bytes \\
    \textbf{Delta (agnostic-layer overhead)} & \textbf{905 bytes} \\
    \hline
  \end{tabular}
\end{table}

\section{Discussion}
\label{sec:eval-discussion}

The results support six conclusions.

First, the cost of the unification is small per call and linear, and its scaling is well
characterised: a cost that grows in direct proportion to both the number of calls and the
registry size in the naive form, dominated by the fixed cost of reopening the
registry file until $R \approx 320$. Second, the single-read optimization eliminates that
call-times-registry-size product, reducing resolution to a cost proportional only to the sum of
the number of calls and the registry size, with a $\sim$200$\times$ speedup in the worst case
(P7, P8, P9) --- a full \emph{identify (P3/P6) $\rightarrow$ fix $\rightarrow$ quantify (P7)} loop.
Third, this optimization has since been propagated to the production resolvers (P10, P11), which
now show the same dominated-by-$N$, flat-in-$F$ shape as the optimized reference path rather than
the call-times-registry-size cost they used to pay.
Fourth, the AWS agnostic packaging layer adds negligible size (S1). Fifth, resolution cost tracks
the number of internal calls rather than the workflow's total call count: external calls add only
the shared tree-traversal cost, never a registry lookup (P15, P16). Sixth, neither nesting depth
nor parallel width changes the cost beyond what the total count of internal calls already
predicts: the resolver's explicit recursion into \texttt{Iteration} and \texttt{Parallel} contexts
is free in practice (P18, P19).

\paragraph{Global framing.}
All measured values lie between microseconds and, in the worst non-optimized write case,
a few hundred milliseconds. Endpoint resolution is therefore not the bottleneck of an integrated
deployment: the dominant cost is the actual cloud deployment (seconds). With the single-read
optimization now applied uniformly across both production resolvers, the overhead of the
unification is asymptotically linear and practically irrelevant in real operation.

\section{Threats to Validity}
\label{sec:eval-threats}

The absolute numbers were obtained on a shared machine with a short JMH configuration; some
experiments used a single fork, which widens error margins. The measurements are consistent among
themselves and with the expected asymptotic forms, but should be reproduced with \texttt{-f 3} on
dedicated hardware before being read as definitive. The evaluation is also intentionally
local-only: it quantifies the resolution and packaging overhead, not the end-to-end cloud
deployment latency, which depends on provider-side behaviour outside the scope of this work.
